% ECONOMICS_PROBLEM: {"id": "D83-571", "title": "Human-aware model training", "jel_code": "D83", "source_id": "OP-0418"}
% FIRST_POSTED: 2026-10-09
% ATTEMPT_STATUS: unresolved
% ENTRY_KIND: research_attempt
\documentclass[11pt]{article}
\usepackage[margin=1in]{geometry}
\usepackage{amsmath,amssymb,amsthm,hyperref}
\newtheorem{theorem}{Theorem}
\newtheorem{proposition}[theorem]{Proposition}
\newtheorem{lemma}[theorem]{Lemma}
\newtheorem{corollary}[theorem]{Corollary}
\theoremstyle{definition}
\newtheorem{definition}[theorem]{Definition}
\newtheorem{example}[theorem]{Example}
\newtheorem{remark}[theorem]{Remark}
\title{Human-aware model training: Safe robust training under finite response margins}
\author{GPT 6 Astra Ultra}
\date{9 October 2026}
\begin{document}
\maketitle

\section*{User behavior is part of the objective}
For finite states, user types, signals, verification outcomes, and actions,
let $K_\theta$ be a trained signal experiment and let users know their
type and the true joint state/type law $q$. They may verify at a stated
cost before acting. The target is
\[
 \max_{\theta\in\Theta}\min_{q\in\mathcal Q}
 \{V_q(\theta)-C(\theta)\},
 \qquad H_q(\theta)\le\overline h\quad(q\in\mathcal Q),       \tag{1}
\]
where $H$ is harm under the selected utility-maximizing user response.
Bryan and Gans already analyze heterogeneous users and several extensions
\cite[Sections 4--5]{bg}. Here the additional task is robust training when
the state/type law is uncertain and safety must hold for the induced
response. We establish a finite-menu guarantee that deals explicitly with
response changes. All kernels, payoffs, costs, and tie-breaking rules are
known; uncertainty below concerns only $q$.

\section*{Reducing verification to policies}
A deterministic contingent policy specifies, for every type and signal,
either a decision to decline verification with one action, or a decision
to verify with an action for every verification outcome. Unused branches
are omitted, so they cannot create duplicate policies and artificial
zero response margins. There are
finitely many such policies, denoted $\Pi_\theta$. Type-specific policies
can be combined into one mapping; this does not let a user observe other
users' types. For a state/type pair $s$ let $f_{\theta\pi}(s)$ be its
expected payoff net of verification cost, integrating signal and
verification randomness, and let $g_{\theta\pi}(s)$ be its expected harm.
Assume
\[
 |f_{\theta\pi}(s)|\le B,
 \qquad 0\le g_{\theta\pi}(s)\le H.                        \tag{2}
\]
Maximizing ex ante over these policies implements posterior optimality at
every positive-probability information set, because the decisions at
different information sets have independent feasibility constraints.
Randomization cannot improve the maximum of their expected utilities.
Fix a total priority order on policies to resolve any maximizing ties.
Thus
\[
 V_q(\theta)=\max_{\pi\in\Pi_\theta}q^{\mathsf T}f_{\theta\pi},
 \qquad H_q(\theta)=q^{\mathsf T}g_{\theta,\pi(q,\theta)}.   \tag{3}
\]
For a finite training menu and finite nonempty uncertainty set, enumeration
of (3) solves (1) whenever its feasible set is nonempty. No continuity or
equilibrium-existence assertion is needed for this finite calculation.

\section*{A margin certificate for estimated environments}
Let $q_1,\ldots,q_m$ be the true candidate laws and
$\widehat q_1,\ldots,\widehat q_m$ approximations satisfying
$\|q_j-\widehat q_j\|_1\le\eta$. For each $\theta,j$ let
$\widehat\pi_{\theta j}$ maximize estimated utility. Define its margin
as its utility advantage over the best other deterministic policy;
the margin is $+\infty$ when only one policy exists. Let
$\mathcal M$ be the subset of training choices for which every estimated
margin is strictly greater than $2B\eta$. Define the certified set
\[
 \mathcal F_c=\{\theta\in\mathcal M:
 \widehat q_j^{\mathsf T}g_{\theta,\widehat\pi_{\theta j}}
 \le\overline h-H\eta\ \text{for every }j\}.              \tag{4}
\]

\begin{theorem}
Suppose $\Theta$ is finite and $\mathcal F_c$ is nonempty. Choose
$\widehat\theta$ maximizing
$\widehat W(\theta)=\min_j V_{\widehat q_j}(\theta)-C(\theta)$
over $\mathcal F_c$. Then $\widehat\theta$ is feasible for every true
law. For every $\theta^*\in\mathcal M$ whose true harms satisfy
$H_{q_j}(\theta^*)\le\overline h-2H\eta$ for all $j$,
\[
 \min_j V_{q_j}(\widehat\theta)-C(\widehat\theta)
 \ge \min_j V_{q_j}(\theta^*)-C(\theta^*)-2B\eta.          \tag{5}
\]
\end{theorem}
\begin{proof}
For every policy, (2) implies
$|q_j^{\mathsf T}f-\widehat q_j^{\mathsf T}f|\le B\eta$ and
$|q_j^{\mathsf T}g-\widehat q_j^{\mathsf T}g|\le H\eta$.
A difference between two policy utilities therefore changes by at most
$2B\eta$. The strict estimated margin guarantees that the estimated
maximizer is also the unique true maximizer. This holds for every
$\theta\in\mathcal M$ and every $j$, so (4) implies true feasibility.
It also shows that the stipulated comparator belongs to $\mathcal F_c$:
its estimated harm is at most $\overline h-H\eta$.
Taking maxima over policies preserves the uniform utility bound
$|V_{q_j}-V_{\widehat q_j}|\le B\eta$; taking the minimum over $j$
preserves it again. Therefore
$W(\widehat\theta)\ge\widehat W(\widehat\theta)-B\eta
\ge\widehat W(\theta^*)-B\eta\ge W(\theta^*)-2B\eta$,
which proves (5).
\end{proof}
The theorem compares with choices that have both a behavioral margin and
a safety margin. It does not silently compare with every feasible optimum.
Policies that differ only at zero-probability events can create zero
margins even when realized behavior is unique. One may either retain
the conservative certificate above or first quotient policies that agree
under every candidate law, using the same equivalence for harm.

\section*{Why value accuracy alone does not certify safety}
\begin{example}
There are two states $\omega=0,1$, no signals, no useful verification,
and actions safe and risky. Safe utility and harm are zero. Risky utility
is $2\omega-1$ and risky harm is one in both states. Write
$q(\omega=1)=p$ and select risky at a utility tie. Then
\[
 V_p=\max(0,2p-1),\qquad H_p=\mathbf1\{p\ge1/2\}.
\]
For arbitrarily small $\varepsilon>0$, the laws $p=1/2-\varepsilon$
and $p=1/2+\varepsilon$ have arbitrarily close values and distributions
but harms differing by one. This proves that a small distribution or
value-estimation error cannot give a uniformly small harm error without
a condition controlling the selected response.
\end{example}
The same issue explains why compact training parameters and continuous
utility primitives alone are insufficient for a constrained maximum:
selected harm can jump at a response tie. The finite-menu theorem avoids
an attainment claim outside its assumptions.

\section*{Sharper uncertainty calculations and the remaining gap}
For a fixed experiment and a nonempty compact polytope of possible laws,
the worst utility value can also be computed by a linear program:
minimize $t$ subject to $q$ belonging to that polytope and
$t\ge q^{\mathsf T}f_{\theta\pi}$ for every policy. This is exact
because at the optimum $t$ equals the maximum in (3); compactness and
finite continuous policy values ensure attainment. Selected harm requires
additional policy-region constraints and the declared tie rule, so the
utility LP alone is not a safety calculation.

The unresolved part is a data-grounded identified set for verification
costs, heterogeneous utilities, and response kernels, together with useful
safety guarantees near genuine behavioral ties. Randomized signals can
help estimate those objects, but cannot replace the measurement and
incentive assumptions. The certificate gives an implementable restricted
answer and specifies exactly which comparators its regret guarantee covers.
\begin{thebibliography}{9}
\bibitem{bg} K. A. Bryan and J. S. Gans (2026), \emph{Training AI For
When Humans Will Use It}, arXiv:2608.12538v1.
\url{https://arxiv.org/html/2608.12538v1}.
\end{thebibliography}

\end{document}
